\section{Support restrictions and the geometry of limiting distances}
\label{shape:section}

Throughout this section we retain the sequence of parameters, scales, and
measures constructed in \cref{scale:selection,scale:limits}.  In particular,
all limiting arrays belong to $X_0$, and all measures are considered on
$\mathcal U=X\setminus\mathcal A$.  As before, $\mathcal D$ denotes the
dyadic rationals.  The constants $\delta$ and $c_b$ are those fixed in
\cref{can:theorem,bridge:theorem}; neither depends on the scales or the
parameter sequence.

For $d\in X_0$, let $P_d$ be the set of times having no open neighborhood
on which $d(x,y)=w|x-y|$ for some constant $w$, and let $Q_d$ be the set
of times having no open neighborhood on which $d(x,y)=|x-y|$.
These sets are closed, and $P_d\subset Q_d$.  Here and below a property of
the distance on an interval concerns every pair of points in that interval.

\begin{proposition}[Support exclusions]\label{shape:exclusions}
There is a measurable set of full measure for $\mu$, $\nu$, and
$\mu_\infty$ on which the following statements hold for every real choice
of their arguments.
\begin{enumerate}
\item There do not exist $r\in\R$ and $t>0$ such that
\begin{equation}\label{shape:forbidden-triple}
 d(r-t,r)>(1-\delta)t,\qquad
 d(r,r+t)>(1-\delta)t,\qquad
 d(r-t,r+t)<\delta t.
\end{equation}
\item If $s\in P_d$ and $t\in Q_d$, then
\begin{equation}\label{shape:slow-pair}
 d(s,t)\le c_b|t-s|.
\end{equation}
\end{enumerate}
The same statements hold almost everywhere for every positive measure
appearing in the decomposition of \cref{scale:decomposition}.
\end{proposition}

\begin{proof}
First fix $r,t\in\mathcal D$, $t>0$.  For all sufficiently large $K$,
$nr$ and $nt$ are integers at every scale $n\ge n_0$ in our sums.
The event \cref{shape:forbidden-triple} for $g_n$ is contained in the
cancellation event of \cref{can:theorem}, at center $nr$ and spacing
$nt$.  Consequently its $\mu_K$-mass is at most
\[
 \frac{C}{\epsilon}\sum_{n_0\le n<N}^{\rm dyadic}
 e^{-\gamma tn}
 \le \frac{C_t}{\epsilon}e^{-\gamma t n_0}\longrightarrow0.
\]
The last conclusion follows from $\epsilon n_0\to\infty$ in
\cref{scale:selection}.  The event is open in $X$.  To pass to a locally
finite limit, test its indicator from below by nonnegative continuous
functions with compact support in $\mathcal U$.  This shows that its
$\mu$-mass is zero.  Taking the countable union over $r,t\in\mathcal D$
and then using continuity of $d$ excludes every real triple: strict
inequalities at a real center and spacing persist at nearby dyadic ones.

We record explicitly the weighted consequence of the bridge estimate
needed for the second exclusion.  Fix a bridge $[u,v]$ with dyadic
endpoints and length $r=v-u>0$.  Fix a power-of-two interval $I$ in the
allowed window about one endpoint and another such interval $I'$ in the
allowed window about the other endpoint.  The two endpoints can be
interchanged.  Put
\[
 W=J(I)+V(I),\qquad H=h(I'),
\]
and fix a threshold $\eta>0$.  These are fixed bounded Lipschitz
functions of finitely many shortfalls, and $W\ge0$.  All their endpoints
align with the integer times at every scale in the sum once $K$ is large.
By \cref{scale:individual-test-bound} and Markov's inequality,
\[
 \sup_{n_0\le n<N}^{\rm dyadic}\PP(H(g_n)>\eta/2)
 \le C_{I',\eta}\epsilon.
\]
Applying \cref{bridge:theorem} at the integer bridge length $rn$, summing,
and using \cref{scale:local-mass}, we obtain
\begin{align}
 &\int \one_{\{d(u,v)>c_b r,\ H>\eta\}}W\,\dd\mu_K
 \nonumber\\
 &\qquad\le C\epsilon\left[
       \int W\,\dd\mu_K+
       \frac1\epsilon\sum_{n_0\le n<N}^{\rm dyadic}
                              M^{-\sigma rn}\right]
 \longrightarrow0.\label{shape:weighted-bridge}
\end{align}
Indeed the first bracketed term is bounded, and the second tends to zero:
the exponential sum is at most $2M^{-\sigma r n_0}$ for all sufficiently
large $K$, whereas $1/\epsilon=o(n_0)$.  All constants and thresholds in
this application are fixed before $K$ tends to infinity.  The lower bounds
on $M$ and on bridge length in \cref{bridge:theorem} therefore hold
eventually; no uniformity over varying test intervals is being assumed.

Suppose now that \cref{shape:slow-pair} fails for some $s\in P_d$,
$t\in Q_d$.  Necessarily $s\ne t$.  Choose nearby distinct dyadic bridge
endpoints so that their distance still exceeds $c_b$ times their temporal
separation.  The endpoints can be chosen close enough to $s,t$ that
sufficiently small neighborhoods of these two times fit in the respective
bridge windows.  Some test $J(I)+V(I)$ with power-of-two length and dyadic
endpoints contained in the first neighborhood is positive.  To see this,
choose such an interval containing $s$ in its interior.  If the tests for
it and all its dyadic subdivisions vanished, successive subdivision,
equal half speeds, and equality in the triangle inequality would make
$d$ affine on that interval, contrary to $s\in P_d$.  In the second
neighborhood there is similarly a power-of-two interval $I'$ with
$h(I')>0$.  In fact exact unit distance between the endpoints of any
interval forces exact unit distance on every subinterval, by the triangle
inequality and $d(x,y)\le|x-y|$; hence an interval containing $t$ in its
interior cannot have zero shortfall.

Choose a positive rational $\eta<h(I')$.  The resulting strict bridge
and shortfall conditions define an open set.  By
\cref{shape:weighted-bridge}, its $W\mu$-mass is zero, using compactly
supported continuous tests in $\mathcal U$ as above.  Its part with $W>0$
therefore has zero $\mu$-mass.  There are only countably many choices of
the bridge, the two tests, their orientation, and the rational threshold.
Every violation is covered by one of these choices, proving
\cref{shape:slow-pair} simultaneously for all real $s,t$.

Both properties are invariant under every translation and positive
dilation of time, including the corresponding division of distances by
the dilation factor.  The measure inequality
$(S^{-1})_*\nu\le\mu$ from \cref{scale:decomposition} consequently transfers
them to $\nu$.  All the other component measures are dominated by $\mu$,
so they inherit them as well.  This also provides a measurable full-measure
set: the exceptional sets were covered by countably many measurable
strict tests before extension to real arguments.
\end{proof}

The local-to-global propagation below uses this zero-hyperbolic
Gromov-product framework \cite[Section~2.2]{Bestvina1997}, with an
explicit argument on the semimetric arrays rather than an invocation
of a theorem about geodesic real trees.

\begin{proposition}[Shape classification]\label{shape:classification}
Every nonaffine array in $X_0$ satisfying \cref{shape:exclusions} has one
of the following forms:
\begin{enumerate}
\item every secant has speed at most $c_b$;
\item there is one exterior ray of exact unit speed, and the complementary
closed half-line has all secant speeds at most $c_b$;
\item there are two exterior rays of exact unit speed, separated by a
closed interval $[a,b]$ of positive length, and every secant contained in
$[a,b]$ has speed at most $c_b$.
\end{enumerate}
The second possibility includes either orientation.  These descriptions
do not assert additivity of distances across the slow part.
The three temporal configurations are illustrated in \cref{shape:figure}.
\end{proposition}

\begin{proof}
We first show that a connected component $U$ of $\R\setminus P_d$ is
affine with a single speed throughout.  Its local speeds agree on
overlapping neighborhoods, and hence are a common $w\in[0,1]$ by
connectedness.  A compact subinterval of $U$ has a sufficiently fine
partition $x_0<\cdots<x_m$ such that each consecutive triple lies in a
local affine neighborhood.  When $w=0$, the triangle inequality
immediately gives $d(x_0,x_m)=0$.  When $w>0$, propagate additivity
inductively.  Suppose it holds from $x_0$ through $x_j$, and use the
four-point inequality of \cref{scale:tree-inequality} at base $v=x_j$ in its
limiting form
\begin{equation}\label{shape:tree-inequality}
 B_v(x,z)\ge\min\{B_v(x,y),B_v(y,z)\},\qquad
 B_v(x,z)=d(v,x)+d(v,z)-d(x,z).
\end{equation}
The inductive additivity gives
$B_v(x_0,x_{j-1})=2d(x_{j-1},x_j)>0$, whereas local additivity gives
$B_v(x_{j-1},x_{j+1})=0$.  Apply
\cref{shape:tree-inequality} to the latter pair with intermediate point
$x_0$.  Nonnegativity of $B_v$ forces $B_v(x_0,x_{j+1})=0$, exactly the
next additivity statement.  The first two steps are already locally
affine.  Thus $d(x_0,x_m)=w(x_m-x_0)$, and varying the compact subinterval
proves the claim.  In particular $P_d\ne\varnothing$ for a nonaffine
array.

Suppose a component $U$ has speed $w>c_b$.  It has at least one finite
endpoint $s\in P_d$, since $U=\R$ would imply affinity.  If $w<1$, every
interior point $t\in U$ belongs to $Q_d$.  Continuity at the endpoint gives
$d(s,t)=w|t-s|$, contradicting \cref{shape:slow-pair}.  Therefore $w=1$.
Such a component cannot have two finite endpoints: they would both belong
to $P_d\subset Q_d$, and their unit mutual speed would again contradict
\cref{shape:slow-pair}.  Hence every component of speed greater than
$c_b$ is an exterior unit ray; there are at most two.

Let $I$ be the closed complement of these rays.  It is the whole line, a
closed half-line, or a closed bounded interval.  For $x<y$ in $I$, if
$[x,y]\cap P_d$ is empty, the two points lie in a component of speed at
most $c_b$.  Otherwise let $p$ and $q$ be the first and last points of
$[x,y]\cap P_d$.  The portions from $x$ to $p$ and from $q$ to $y$ have
speed at most $c_b$, by the component claim and continuity.  The middle
portion satisfies the same bound by \cref{shape:slow-pair}, because
$p,q\in P_d\subset Q_d$.  The triangle inequality gives
$d(x,y)\le c_b(y-x)$.  The only case not yet settled is a bounded slow
interval of length zero; \cref{shape:crossloss} below excludes it.
\end{proof}

\begin{lemma}[Loss across two rays]\label{shape:crossloss}
Suppose there are two exterior unit rays with slow interval $[a,b]$,
and put $L=b-a$.  Then $L>0$, and, with the fixed constant
\begin{equation}\label{shape:cross-constant}
 C_*=2+\frac6\delta,
\end{equation}
every $z\in\R$ and $t>0$ satisfy
\begin{equation}\label{shape:cross-bound}
 0\le 1-\frac{d(z,z+t)}t\le\frac{C_*L}{t}.
\end{equation}
The shortfall is zero whenever the interval is wholly in either closed
unit ray.
\end{lemma}

\begin{proof}
For the moment permit $L=0$.  Write
$x_u=a-u$, $y_v=b+v$, where $u,v\ge0$, and put $D=d(a,b)\le L$.
Set
\[
 s=\frac{u+v-D-d(x_u,y_v)}2.
\]
If $s\le0$, the absolute shortfall of $[x_u,y_v]$ is at most
$L+D\le2L$.  Suppose $s>0$.  The reverse triangle inequalities imply
$d(x_u,y_v)\ge u-D-v$ and $d(x_u,y_v)\ge v-D-u$, so that
$s\le\min\{u,v\}$.  At base $a$ the unit rays and triangle inequalities
give
\begin{align*}
 B_a(x_u,x_s)&=2s,\\
 B_a(x_u,y_v)&\ge u+(v-D)-d(x_u,y_v)=2s,\\
 B_a(y_v,y_s)&\ge(v-D)+(s-D)-(v-s)=2s-2D.
\end{align*}
Twice applying \cref{shape:tree-inequality} yields
$B_a(x_s,y_s)\ge2s-2D$.  As $d(a,x_s)=s$ and
$d(a,y_s)\le D+s$, it follows that
\begin{equation}\label{shape:matched-ray-distance}
 d(x_s,y_s)\le3D\le3L.
\end{equation}

Use the temporal midpoint $m=(a+b)/2$ and spacing $t_*=s+L/2$.
Its two wing distances obey
\[
 d(x_s,m),\ d(m,y_s)\ge s-L/2.
\]
If $L>0$ and $s>3L/\delta$, these bounds are strictly greater than
$(1-\delta)t_*$ and \cref{shape:matched-ray-distance} is strictly less
than $\delta t_*$.  This is the forbidden triple of
\cref{shape:forbidden-triple}.  If $L=0$, every $s>0$ produces that same
strict contradiction.  Thus always $s\le3L/\delta$.  A full crossing
interval has absolute shortfall
\[
 (u+L+v)-d(x_u,y_v)=L+D+2s\le C_*L.
\]
An interval meeting only part of the slow interval has absolute
shortfall at most twice the length of that part.  For example, if
$x<a\le y\le b$, then
$d(x,y)\ge(a-x)-d(a,y)\ge(a-x)-(y-a)$, giving loss at most $2(y-a)$.
The other orientation is identical.  An interval inside $[a,b]$ has loss
at most its length, hence at most $L$, and intervals in the unit rays
have zero loss.  This proves \cref{shape:cross-bound} in every case.
If $L=0$, these conclusions would give exact unit distance between every
pair of times, an affine array.  Nonaffinity therefore forces $L>0$.
\end{proof}

\begin{figure}[tb]
\centering
\begin{tikzpicture}[x=1.05cm,y=1cm,>=Latex,font=\small]
  \node[anchor=east] at (-4.4,2.7) {No unit ray};
  \draw[gray!65,line width=3pt] (-4,2.7)--(4,2.7);
  \node[fill=white,inner sep=2pt] at (0,2.7) {all secants have speed $\le c_b$};
  \node[anchor=east] at (-4.4,1.5) {One unit ray};
  \draw[blue!65!black,line width=1.1pt,->] (-1,1.5)--(-4,1.5);
  \draw[gray!65,line width=3pt] (-1,1.5)--(4,1.5);
  \fill (-1,1.5) circle (1.6pt);
  \node[above] at (-2.6,1.5) {unit speed};
  \node[above] at (1.5,1.5) {slow half-line};
  \node[below] at (-1,1.5) {$a$};
  \node[anchor=east] at (-4.4,0) {Two unit rays};
  \draw[blue!65!black,line width=1.1pt,->] (-1,0)--(-4,0);
  \draw[gray!65,line width=3pt] (-1,0)--(1,0);
  \draw[blue!65!black,line width=1.1pt,->] (1,0)--(4,0);
  \draw[<->] (-1,.55)--(1,.55) node[midway,above] {$L=b-a>0$};
  \foreach \x in {-3.3,-1,0,1,3.3}{\fill (\x,0) circle (1.6pt);}
  \node[below] at (-3.3,0) {$x_s=a-s$};
  \node[below] at (-1,0) {$a$};
  \node[below] at (0,0) {$m$};
  \node[below] at (1,0) {$b$};
  \node[below] at (3.3,0) {$y_s=b+s$};
\end{tikzpicture}
\caption{The possible temporal regions in \cref{shape:classification};
the one-ray picture may also be reflected.  A slow region need not be
affine.  In the bottom picture, excessive loss between the exterior rays
would produce the symmetric forbidden triple $(x_s,m,y_s)$ used in
\cref{shape:crossloss}.  The diagram describes time regions, not an
embedding of the semimetric.}
\label{shape:figure}
\end{figure}

\FloatBarrier
\section{A capped deficit and the final balance}
\label{end:section}

The cap will compare the slow interval of a limiting shape at fine and
coarse scales. It must equal one on uniformly slow secants. For a
two-ray shape, \cref{shape:cross-bound} makes the shortfall of a long
crossing interval small; a sufficiently small slope of the cap near
zero will make its averaged coarse contribution smaller than the
fine-scale occupancy of the slow interval. This is an integrated
comparison, not a pointwise sign assertion for the observable below.

Choose a fixed $\alpha>0$ so small that $10C_*\alpha<1$.  Choose a smooth
nondecreasing function $\Phi:[0,1]\to[0,1]$ such that
\begin{equation}\label{end:cap}
 \Phi(u)=\alpha u\ \text{near }0,\qquad
 \Phi(u)=1\quad(1-c_b\le u\le1).
\end{equation}
For completeness, one can choose a nonnegative smooth derivative equal
to $\alpha$ near zero, zero near $1-c_b$, and with integral one on
$[0,1-c_b]$, by adding a nonnegative bump supported strictly between the
two endpoint neighborhoods.  Integrating and continuing constantly gives
\cref{end:cap}.  Positivity near zero and compactness give constants
$0<c_\Phi\le C_\Phi<\infty$ with
\begin{equation}\label{end:linear-bounds}
 c_\Phi u\le\Phi(u)\le C_\Phi u\qquad(0\le u\le1).
\end{equation}
Extend $\Phi$ smoothly to the real line, keeping it linear immediately
to the left of zero.  The extension only defines continuous tests on
$X$; the arrays used below have the relevant shortfalls in $[0,1]$.

Write $h_t(z)=1-d(z,z+t)/t$, and abbreviate $J=J([0,2])$ and
$V=V([0,2])$.  Define
\begin{equation}\label{end:G-definition}
 G(d)=\Phi(h_2(0))-\tfrac12\Phi(h_1(0))-\tfrac12\Phi(h_1(1)).
\end{equation}
If $a=h_1(0)$, $b=h_1(1)$, and $\bar a=(a+b)/2$, then
$h_2(0)=\bar a+J$ and $V=(a-b)^2$.  The mean value theorem followed by
the second-order symmetric Taylor bound yields
\begin{equation}\label{end:G-bound}
 |G|\le\|\Phi'\|J+\frac18\|\Phi''\|V\le C(J+V),
\end{equation}
where the derivative norms can be taken on the compact argument range
of these three tests on $X$.  In particular $G$ is $\mu$-integrable by
\cref{scale:limiting-test-bound}.

At the terminal scale, $\int h_1(0)\,\dd\nu_K=1$.  After a subsequence,
\cref{end:linear-bounds} therefore gives a limit
\begin{equation}\label{end:H-definition}
H_\Phi=\lim_K\int\Phi(h_1(0))\,\dd\nu_K,
 \qquad c_\Phi\le H_\Phi\le C_\Phi.
\end{equation}

There are three steps to the contradiction. First, dyadic telescoping
will give $\int G\,\dd\mu=H_\Phi$, with no loss of weighted mass to
$\mathcal A$. Then the positive decomposition of
\cref{scale:decomposition} separates this integral into a terminal
contribution and the contribution of the invariant remainder. We prove
that the former is strictly less than $H_\Phi$ and the latter is
nonpositive. The cap's two prescribed behaviors enter precisely in
these comparisons.

\begin{proposition}[Exact limiting balance]\label{end:balance}
With the preceding choices,
\begin{equation}\label{end:balance-equation}
 \int G\,\dd\mu=H_\Phi>0.
\end{equation}
\end{proposition}

\begin{proof}
Stationarity makes the two half-interval terms in
\cref{end:G-definition} have equal expectation.  Since
$h_2(0;g_n)=h_1(0;g_{2n})$, summing over the dyadic scales gives exactly
\[
 \int G\,\dd\mu_K
 =\int\Phi(h_1(0))\,\dd\nu_K
  -\frac1\epsilon\E\Phi(h_1(0;g_{n_0})).
\]
The final term tends to zero by
\cref{scale:initial-test-bound,end:linear-bounds}.  Thus these integrals
converge to $H_\Phi$.  We must show that this mass does not escape to
$\mathcal A$ under vague convergence.

The measures $J\mu_K$ and $V\mu_K$ have uniformly bounded total mass on
the compact space $X$, and the signed measures $G\mu_K$ have uniformly
bounded total variation by \cref{end:G-bound}.  Take simultaneous weak
limits
\[
 J\mu_K\longrightarrow\lambda_J,\qquad
 V\mu_K\longrightarrow\lambda_V,\qquad
 G\mu_K\longrightarrow\eta.
\]
On $\mathcal U$ they agree with $J\mu$, $V\mu$, and $G\mu$ respectively:
multiply any compactly supported continuous test there by the indicated
continuous weight.  We show that $\eta|_{\mathcal A}=0$ by treating all
three ranges of affine speeds.  Write $d_w(s,t)=w|s-t|$.

\paragraph{Speeds $0\le w\le c_b$.}
Both $\Phi'$ and $\Phi''$ vanish at every shortfall $1-w$ in this compact
speed range, including its boundary $w=c_b$.  For each $\rho>0$ there is
an open neighborhood $U_\rho$ of
$\{d_w:0\le w\le c_b\}$ on which the Taylor argument for
\cref{end:G-bound} gives
\[
 |G|\le\rho(J+V).
\]
Indeed all arguments between the three relevant shortfalls are then
uniformly close to $[1-c_b,1]$, where both derivatives vanish.
For a continuous test supported in $U_\rho$, pass this domination to
the weak limits.  It implies
$|\eta||_{U_\rho}\le\rho(\lambda_J+\lambda_V)|_{U_\rho}$.
The latter positive measure is finite on $X$.  Letting $\rho$ decrease
to zero shows that $\eta$ has zero restriction to this compact set of
affine arrays.

\paragraph{Speeds $c_b<w<1$.}
Choose a fixed positive dyadic $R$ so large that the endpoint windows of
the bridge $[0,R]$ contain $[0,2]$ and $[R,R+1]$ respectively; for
example require $\kappa R>2$.  Put $W=J+V$ and
$H=h([R,R+1])$.  For each positive rational $\eta_0$, the open set
\[
 U_{\eta_0}=\{d:d(0,R)>c_bR,\ h([R,R+1])>\eta_0\}
\]
has $W\mu_K$-mass tending to zero by
\cref{shape:weighted-bridge}.  The open-set inequality for weak
convergence of positive measures gives
$(\lambda_J+\lambda_V)(U_{\eta_0})=0$.  Each $d_w$ in the current
speed range belongs to such a set by taking $0<\eta_0<1-w$.
This countable cover and \cref{end:G-bound} show that $\eta$ vanishes
on all intermediate affine speeds.  Notice that the threshold is fixed
separately for each test; no threshold uniform as $w\uparrow1$ is needed.

\paragraph{The speed $w=1$.}
Near $d_1$, all three arguments in \cref{end:G-definition} are in the
linear part of $\Phi$, so $G=\alpha J$.  It remains to prove
$\lambda_J(\{d_1\})=0$.  Fix an inverse power of two $q\in(0,1)$.
For the product over $[0,2n]$, with center $n$, put
\[
 D=n-\min\{g(0,n),g(n,2n)\},\qquad
 \ell=g(0,n)+g(n,2n)-g(0,2n)=2nJ(g_n).
\]
The quantities $D/n$ and $\ell/n$ are continuous tests of $g_n$,
both zero at $d_1$.  Thus the condition
$2q-2D/n>\ell/n$ defines an open neighborhood $U_q$ of $d_1$.
Eventually $s=qn$ is an integer at every scale in the sum.
On $U_q$, the truncation estimate of \cref{scale:truncation} gives,
with its fixed additive constant $A$,
\[
 \ell\le\ell_{\rm cut}+A,
 \qquad
 J\le qJ([1-q,1+q])+\frac{A}{2n}.
\]
The shorter-interval test telescopes exactly, by stationarity:
\[
 \int J([1-q,1+q])\,\dd\mu_K
 =\frac{e_{qN}-e_{qn_0}}\epsilon\le1,
\]
where dyadic monotonicity gives $e_{qN}\le e_N=\epsilon$.
Also $\epsilon^{-1}\sum_n n^{-1}=O((\epsilon n_0)^{-1})=o(1)$.
Since $J\ge0$,
\[
 \int_{U_q}J\,\dd\mu_K\le q+o(1).
\]
The open-set inequality now gives $\lambda_J(U_q)\le q$, so
$\lambda_J(\{d_1\})\le q$.  First take the fixed weak limit and then
let $q$ decrease through inverse powers of two to obtain zero.  The
local identity $G=\alpha J$ transfers this conclusion to $\eta$.

These cases cover $\mathcal A$.  Weak convergence on compact $X$
preserves the total mass of the signed measures, so
$\eta(X)=H_\Phi$.  As $\eta|_{\mathcal A}=0$ and
$\eta|_{\mathcal U}=G\mu$, this is \cref{end:balance-equation}.
\end{proof}

\subsection{The terminal measure and its backward dilates}

The decomposition in \cref{scale:decomposition}, together with
\cref{end:G-bound}, gives an absolutely convergent identity
\begin{equation}\label{end:integrated-decomposition}
 \int G\,\dd\mu
 =\int G\,\dd\mu_\infty
  +\sum_{j=1}^{\infty}\int G(S^{-j}d)\,\dd\nu(d).
\end{equation}
More precisely, the sum of the integrals of $|G(S^{-j}d)|$ is at most
$\int|G|\,\dd\mu$, so all the subsequent rearrangements of these
component integrals are justified.

\begin{lemma}[Strict bound for the terminal contribution]
\label{end:terminal}
The series in \cref{end:integrated-decomposition} satisfies
\begin{equation}\label{end:terminal-strict}
 \sum_{j=1}^{\infty}\int G(S^{-j}d)\,\dd\nu(d)<H_\Phi.
\end{equation}
\end{lemma}

\begin{proof}
Nonnegative compact cutoffs exhausting $\mathcal U$, the terminal
prelimit identities, and monotone convergence imply
\begin{equation}\label{end:terminal-bounds}
 \int h_1(0)\,\dd\nu\le1,
 \qquad \int\Phi(h_1(0))\,\dd\nu\le H_\Phi.
\end{equation}
For each integer $m\ge0$, subdivision and the triangle inequality give
\[
 h_1(0)\ge2^{-m}\sum_{r=0}^{2^m-1}h_{2^{-m}}(r2^{-m}).
\]
Dyadic translation invariance of $\nu$ therefore implies
$\int h_{2^{-m}}(0)\,\dd\nu\le1$.  The same shortfalls inserted into
$\Phi$ have finite integrals, bounded by $C_\Phi$.

Set $A_m=\int\Phi(h_{2^{-m}}(0))\,\dd\nu$.  Translation invariance
at the shift $2^{-j}$ gives
\[
 \int G(S^{-j}d)\,\dd\nu=A_{j-1}-A_j.
\]
Absolute convergence of the series implies convergence of $A_m$.
Fatou's lemma yields
\begin{align}
 \sum_{j=1}^{\infty}\int G(S^{-j}d)\,\dd\nu
 &\le A_0-
    \int\liminf_{m\to\infty}\Phi(h_{2^{-m}}(0))\,\dd\nu
 \nonumber\\
 &\le H_\Phi-
    \int\liminf_{m\to\infty}\Phi(h_{2^{-m}}(0))\,\dd\nu.
 \label{end:terminal-fatou}
\end{align}
If $\nu=0$, the series is zero and \cref{end:H-definition} proves the
strict conclusion.  Otherwise every shape in its full-measure support
has a slow interval with nonempty interior by
\cref{shape:classification,shape:crossloss}.  There is a dyadic $z$ in
that interior.  For all sufficiently large $m$, the interval
$[z,z+2^{-m}]$ is slow, so
$\Phi(h_{2^{-m}}(z))=1$ by \cref{end:cap}.

The measurable sets
\[
 E_z=\{d:\Phi(h_{2^{-m}}(z))=1\text{ for every sufficiently large }m\},
 \qquad z\in\mathcal D,
\]
cover this full-measure support.  At least one has positive $\nu$-measure,
and translation invariance gives $\nu(E_0)=\nu(E_z)>0$.
Consequently the subtracted integral in \cref{end:terminal-fatou} is
positive.  It is finite by Fatou's lemma and the bound $A_m\le C_\Phi$.
This proves \cref{end:terminal-strict} in all cases.
\end{proof}

\subsection{The translation- and dilation-invariant remainder}

We next give the measurable and local-finiteness details needed to
integrate over shapes.  A dyadic interval whose endpoint distance equals
its length has exact unit speed throughout, by the triangle inequality
and the Lipschitz bound.  Thus the local unit-speed set is the union of
the interiors of those countably many intervals.  More directly, define
extended-real endpoint functions on $X_0$ by
\begin{align}
 a(d)&=\sup\{r\in\mathcal D:
       d(p,q)=q-p\text{ for all }p,q\in\mathcal D, p<q<r\},
 \label{end:left-endpoint}\\
 b(d)&=\inf\{r\in\mathcal D:
       d(p,q)=q-p\text{ for all }p,q\in\mathcal D, r<p<q\}.
 \label{end:right-endpoint}
\end{align}
Use $\sup\varnothing=-\infty$ and $\inf\varnothing=+\infty$.
Each ray condition is a countable intersection of closed secant
conditions, so $a,b$ are measurable.  On our nonaffine shapes a finite
$a$ is exactly the right endpoint of the exterior left unit ray, and a
finite $b$ is exactly the left endpoint of the exterior right unit ray.
They obey, for $r\in\mathcal D$,
\begin{equation}\label{end:endpoint-covariance}
 a(\tau_r d)=a(d)-r,\quad b(\tau_r d)=b(d)-r,
 \qquad a(Sd)=a(d)/2,\quad b(Sd)=b(d)/2.
\end{equation}
Thus the three shape classes and the length $L=b-a$ on the two-ray
class are measurable and transform as asserted.

\begin{lemma}[Nonpositive invariant contribution]\label{end:remainder}
The invariant remainder satisfies
\begin{equation}\label{end:remainder-bound}
 \int G\,\dd\mu_\infty\le0.
\end{equation}
\end{lemma}

\begin{proof}
On an entirely slow shape, each shortfall in
\cref{end:G-definition} is at least $1-c_b$, so all three capped values
equal one and $G=0$.

Consider next shapes with exactly one unit ray.  First take the
orientation with a left unit ray and finite interface $a$.
For a bounded interval $[A,B]$ of possible interface locations, choose
one fixed dyadic secant strictly to the left of $A$ and a second strictly
to the right of $B$.  Their speed ratios on every such shape are
respectively $1$ and at most $c_b$.  The arrays with these two prescribed
closed conditions form a compact subset of $X_0$ disjoint from
$\mathcal A$: an affine array would have the same speed on both secants.
Local finiteness of $\mu_\infty$ on $\mathcal U$ therefore implies finite
mass for interfaces in $[A,B]$.

Push the restriction of $\mu_\infty$ to this shape class forward by $a$,
obtaining a measure $\lambda$ on $\R$ finite on bounded intervals.
Its invariances, from \cref{end:endpoint-covariance}, give
\[
 \lambda([0,1))=\lambda([0,2))
 =\lambda([0,1))+\lambda([1,2))=2\lambda([0,1)).
\]
The first equality uses $S$-invariance; the last uses translation by
one.  Finiteness forces $\lambda([0,1))=0$, and integer translates
covering $\R$ show that this entire shape class has zero mass.  The same
argument with left and right reversed treats the other one-ray class.
Half-open intervals make the argument valid without any advance
assumption about endpoint atoms.

It remains to consider the two-ray class.  Restrict to the length band
\[
 \mu_*:=\mu_\infty\big|_{\{\text{two rays},\ 1\le L<2\}}.
\]
This measure is invariant under dyadic translations.  Its left endpoint
has finite intensity on bounded intervals, as follows explicitly.
Partition any bounded range of possible $a$ into finitely many bins
$[k/4,(k+1)/4]$.  In each bin the fixed secant
$[(k+1)/4,(k+2)/4]$ lies in $[a,b]$, since $L\ge1$, and hence has ratio
at most $c_b$.  A fixed dyadic secant to the left of the whole bounded
range has ratio one.  The two conditions again put each bin's arrays
inside a compact subset of $\mathcal U$.  A finite union gives the
claimed finite intensity.  The right endpoint also has finite intensity
on bounded intervals because $1\le b-a<2$ bounds $a$ whenever it bounds
$b$.

Neither endpoint intensity has an atom.  If the mass at some point were
positive, dyadic translation invariance would give the same positive
mass to infinitely many distinct points in a bounded interval,
contradicting local finiteness.  Define the finite number
\[
 I_0=\mu_*\{0\le a<1\}.
\]
Every integer interval of length one has this same endpoint mass.

For $j\in\Z$, let $B_j$ be the two-ray band
$2^j\le L<2^{j+1}$.  Since $S$ divides lengths by two and preserves
$\mu_\infty$,
\[
 \mu_\infty|_{B_j}=(S^{-j})_*\mu_*.
\]
The bands are disjoint and cover the two-ray class.  Integrability of
$|G|$ permits reindexing $j$ to $-j$, yielding
\begin{equation}\label{end:band-series}
 \int_{\rm two\ rays}G\,\dd\mu_\infty
 =\sum_{j\in\Z}\int G(S^j d)\,\dd\mu_*(d),
 \qquad
 \sum_{j\in\Z}\int|G(S^j d)|\,\dd\mu_*<\infty.
\end{equation}

Put $F_j=\int\Phi(h_{2^j}(0))\,\dd\mu_*$.  These quantities are
finite.  Indeed \cref{shape:crossloss} makes the integrand zero unless
$[0,2^j]$ meets $[a,b]$, which requires $-2\le a\le2^j$; the endpoint
intensity on this bounded interval is finite.  Translation invariance
at the dyadic shift $2^j$ gives
\begin{equation}\label{end:F-difference}
 \int G(S^j d)\,\dd\mu_*=F_{j+1}-F_j.
\end{equation}
The absolute convergence in \cref{end:band-series} implies that $F_j$
has finite limits both as $j\to-\infty$ and as $j\to+\infty$.

For the fine-scale limit, use the dominating integrable function
$\one_{\{-2\le a\le1\}}$ for $j\le0$.  Off the null sets
$\{a=0\}$ and $\{b=0\}$, if $a<0<b$, then $[0,2^j]$ is contained
in the slow interval for all sufficiently negative $j$, and the capped
shortfall equals one.  If zero lies outside the slow interval, the
small interval is eventually in a unit ray and its shortfall is zero.
Dominated convergence therefore gives
\begin{equation}\label{end:fine-limit}
 F_{\rm fine}:=\lim_{j\to-\infty}F_j
 =\mu_*\{a<0<b\}\ge\mu_*\{-1\le a<0\}=I_0.
\end{equation}
The inequality uses $L\ge1$; possible equality endpoints have zero mass.

For the coarse-scale limit, \cref{shape:cross-bound} gives
$0\le h_t(0)\le2C_*/t$ throughout this band.  Once $t$ is sufficiently
large, all these shortfalls lie in the linear region of \cref{end:cap},
uniformly in the shape.  Hence
\[
 \int\Phi(h_t(0))\,\dd\mu_*
 \le \frac{2\alpha C_*}{t}\,
                  \mu_*\{-2\le a\le t\}.
\]
For the large dyadic integers $t=2^j$, atomlessness and integer
translation invariance make the last mass $(t+2)I_0$ (the weaker bound
$(t+3)I_0$ would also suffice).  Thus
\begin{equation}\label{end:coarse-limit}
 F_{\rm coarse}:=\lim_{j\to+\infty}F_j
 \le2\alpha C_*I_0.
\end{equation}
Summing \cref{end:F-difference} and applying
\cref{end:fine-limit,end:coarse-limit} now proves
\[
 \int_{\rm two\ rays}G\,\dd\mu_\infty
 =F_{\rm coarse}-F_{\rm fine}
 \le-(1-2\alpha C_*)I_0\le0.
\]
This argument also covers $I_0=0$; in that case integer translates of
$\{0\le a<1\}$ show that $\mu_*=0$.  Together with the entirely slow
and one-ray cases, this establishes \cref{end:remainder-bound}.
\end{proof}

\begin{theorem}[Contradiction and eventual positive entropy]
\label{end:contradiction}
There is no sequence $k\to+\infty$ along which the largest Lyapunov
exponent of $f_k$ is zero almost everywhere.  Consequently
\cref{main:entropy} holds.
\end{theorem}

\begin{proof}
Under the supposed zero-exponent sequence, all the preceding measures
and identities were constructed.  Combining
\cref{end:integrated-decomposition,end:terminal-strict,end:remainder-bound}
gives
\[
 \int G\,\dd\mu<H_\Phi,
\]
contradicting the exact balance \cref{end:balance-equation}.

For every parameter, the top Lyapunov exponent exists almost everywhere
and is nonnegative, since the derivative cocycle has determinant one.
If positivity on a set of positive area failed at arbitrarily large
positive parameters, the exponent would therefore be zero almost
everywhere at each member of a sequence tending to $+\infty$, which has
just been excluded.  There is consequently a finite $k_0$ such that,
for every $k>k_0$, the top exponent is positive on a set of positive
normalized area.  Increasing the threshold gives the same statement
for every $k\ge k_0$.  The entropy formula \cref{main:pesin} gives
$h_m(f_k)>0$ for every such $k$.  In particular the measurable parameter
interval $[k_0+1,k_0+2]$ has positive one-dimensional Lebesgue measure
and satisfies the required assertion.
\end{proof}
